\paragraph{Benchmarks}
As program synthesis research advances, the need for rigorous and multidimensional benchmarks has become increasingly important. Most existing coding benchmarks concentrate on functional correctness, often assessed through manually crafted test cases and isolated execution environments. For instance, HumanEval provides expert-written programming tasks paired with curated test suites, while MBPP\cite{austin2021programsynthesislargelanguage} collects problems through crowdsourcing. To improve coverage and scalability, more recent approaches have turned to automated test generation\cite{liu2023codegeneratedchatgptreally}, enabling broader application in areas such as software engineering and developer assistance\cite{jimenez2024swebenchlanguagemodelsresolve}.


In addition to correctness, performance profiling is gaining traction as a crucial axis in benchmarking (\cite{shypula2024learningperformanceimprovingcodeedits};\cite{liu2023codegeneratedchatgptreally};\cite{huang2025effibenchbenchmarkingefficiencyautomatically};\cite{qiu2025efficientllmgeneratedcoderigorous}). However, most existing frameworks emphasize competition-style tasks or sequential execution and provide limited support for parallel or hardware-specific evaluation. While a few benchmarks address parallel CPU programming (\cite{Nichols_2024};\cite{chaturvedi2024hpccoderv2studyingcodellms}), GPU-targeted benchmarks remain scarce, despite the growing importance of hardware efficiency in deploying large-scale deep learning models.

To address this gap, benchmarks like KernelBench~\cite{ouyang2025kernelbenchllmswriteefficient} and TritonBench~\cite{li2025tritonbenchbenchmarkinglargelanguage} have emerged, offering domain-specific assessments of LLM performance on GPU kernel generation. KernelBench~\cite{ouyang2025kernelbenchllmswriteefficient} focuses on correctness and runtime efficiency across diverse workloads using metrics like "fast\_p", while TritonBench~\cite{li2025tritonbenchbenchmarkinglargelanguage} highlights the challenges LLMs face in generating performant code for the Triton DSL. These benchmarks reveal a substantial gap between general-purpose code generation and the demands of high-performance, hardware-aware synthesis, underscoring the need for specialized, compute-oriented evaluation frameworks.

\paragraph{LLMs for Code Generation}

More recent work has explored the use of large language models for code generation, but often in general-purpose contexts with limited success on performance-critical low level code, such as DeepSeek-Coder\cite{guo2024deepseekcoderlargelanguagemodel} and Qwen-Coder\cite{hui2024qwen25codertechnicalreport}, which have achieved strong performance on broad coding benchmarks.

There are also LLMs specifically trained for the kernel generation task, such as KernelLLM. However, its capability is limited to translating PyTorch code into Triton kernels, restricting its general applicability. In contrast, our agent can generate Triton code based on both natural language instructions and reference code, allowing for more flexible and context-aware code synthesis. Furthermore, specialized models like KernelLLM often lack the broader reasoning and problem-solving capabilities exhibited by larger state-of-the-art LLMs such as GPT-4.1 and Claude 4, which are critical for tackling complex optimization tasks involving multiple constraints and trade-offs.


\paragraph{Inference-Time Compute Scaling}
Inference-time compute scaling refers to the strategy of allocating increased computational resources—such as longer context windows, multiple reasoning passes, or parallel generations — during the inference phase of machine learning models to boost performance on complex tasks.
Notably, this approach enables significant improvements in large language model (LLM) performance without requiring additional training.
For example, the Chain-of-Thought~\cite{shinn2023reflexionlanguageagentsverbal} prompting technique enhances reasoning by encouraging models to answer questions in a step-by-step manner. Similarly, the Reflexion~\cite{shinn2023reflexionlanguageagentsverbal} framework improves agent performance by integrating verbalized feedback into an episodic memory buffer, which guides more informed decision-making in future attempts.
Our proposed GEAK agent also leverages inference-time compute scaling to improve code generation and optimization quality. In addition, we incorporate a Reflexion-style module for debugging, enabling the agent to iteratively refine its outputs based on past execution feedback.
